\begin{table}[htbp]
  \centering
  \caption{Decomposition of Sectoral Labor Productivity Growth by Subperiod}
  \label{tab:ex1_decomp_sub}
  \begin{threeparttable}
  \centering
  \begin{tabular}{l *{6}{S[table-format=-1.2]}}
    \toprule
     & \multicolumn{3}{c}{Panel A. 1950--1980} & \multicolumn{3}{c}{Panel B. 1980--2023} \\
    \cmidrule(lr){2-4} \cmidrule(lr){5-7}
     & {Goods} & {RI-Serv.} & {NRI-Serv.} & {Goods} & {RI-Serv.} & {NRI-Serv.} \\
    \midrule
    I.\ Endogenous innovation & 1.88 & 0.62 & 0.59 & 2.14 & 0.94 & 0.78 \\
    II.\ Embodied technical progress & 0.35 & 1.00 & 0.45 & 0.35 & 1.00 & 0.45 \\
    III.\ Relative price effect & -0.02 & 0.36 & 0.12 & -0.24 & -0.03 & 0.15 \\
    IV.\ Labor efficiency & 0.37 & 0.34 & 0.35 & 0.37 & 0.34 & 0.34 \\
    \midrule
    Total & 2.58 & 2.32 & 1.51 & 2.62 & 2.25 & 1.72 \\
    \bottomrule
  \end{tabular}
  \begin{tablenotes}[flushleft]
    \footnotesize
    \RaggedRight
    \item[] \textit{Notes:} Entries report annualized contributions, in percentage points, to sectoral labor productivity growth over 1950--1980 (Panel A) and 1980--2023 (Panel B), where sectoral labor productivity is $\mathcal{A}_j = Y_j/L_j^p$. The decomposition follows Table 3. Entries are rounded so that they add to the printed totals.
  \end{tablenotes}
  \end{threeparttable}
\end{table}
